\section{Introduction}{\label{tab:intro}}

Word embedding models have been designed for representing the meaning of words in a vector space. These models have become a fundamental NLP technique and have been widely used in various applications. However, prior studies show that such models learned from human-generated corpora are often prone to exhibit social biases, such as gender stereotypes~\cite{kaiwei,caliskan2017semantics}. 
For example,  the word ``programmer'' is neutral to gender by its definition, but an embedding model trained on a news corpus  associates ``programmer'' closer with ``male'' than  ``female''. 

Such a bias substantially affects downstream applications.  \newcite{zhao2018gender} show that a coreference resolution system is sexist due to the word embedding component used in the system. This concerns the practitioners who use the embedding model to build gender-sensitive applications such as a resume filtering system or a job recommendation system as the automated system may discriminate candidates based on their gender, as reflected by their name.
Besides, biased embeddings may implicitly affect downstream applications used in our daily lives. For example, when searching for ``computer scientist'' using a search engine, as this phrase is closer to male names than female names in the embedding space, a search algorithm using an embedding model in the backbone tends to rank male scientists higher than females', hindering women from being recognized and further exacerbating the gender inequality in the community.

\iffalse
Word embedding, which assigns each word a low dimensional representation that preserves semantic proximity,  has become a fundamental NLP technique.
However, word embedding models that are learned from human-generated corpora which often exhibit social biases, such as gender stereotypes~\cite{kaiwei}. 
Consequently, the word embeddings may carry incorrect gender information, even for  gender-neutral words. For example, ``doctor'' has a higher proximity to ``male'' than to ``female'' in the embedding space while ``nurse'' is opposite.
Such bias substantially affects downstream applications, as the inappropriate proximity of ``doctor'' and ``nurse'' can lead to biased job recommendations because of the gender information of job seekers registered in the system.

%Also as \newcite{caliskan2017semantics} shows that comparing with art terms ``dance'' and ``poetry'', male terms such as ``man'', ``boy'' are more associated with science terms ``math'', ``algebra'', while female terms like ``woman'', ``girl'' are more associated with the art.
In addition, current search engines utilize word embedding technique. In the embedding space, ``computer science'' is closer to male names than female~\cite{kaiwei}. When a user searches ``CMU Computer Science PhD student'', between two pages that have same framework and differ only in the names, the searching algorithm would rank male's web page higher than the female's. 
Such gender stereotype in word embeddings hinders women from being recognized as computer scientists and will worsen the gender inequality in the computer society.
\fi
% For example, nbg~\cite{kaiwei}.  
% This will cause the system to rank a male page higher than a female one even though the only difference is the name which in result will widen the gap between different genders in computer science.
%Because of the amplification of gender bias in the word embedding training process, ``Computer Science'' in the continuous word space is much closer to male names like John rather than female names like Marry, leading to the unpleasant search experiences. 

%The existence of gender bias has been observed and documented in literature. 
%\newcite{caliskan2017semantics} argues that applying machine learning algorithms to ordinary human language may result in human-like semantic biases. 
% For example, words in the category of \textit{Math} are found closer to \textit{male} than to \textit{female} in the latent representation space, while \textit{Arts}-related words turn out just the opposite. 


%, among which gender bias is drawing increasing attentions recently~\cite{jieyu}. 

%An example  exists in the fallacy that ``nursing'' is a woman biased job but a man is more likely to be a physician only because of their genders. Word embeddings trained upon \textit{GoogleNews} strongly support this fallacy, since we find the cosine similarity between the word vectors of ``nurse'' and ``he'' is much smaller than that between ``nurse'' and ``she''. The same bias also exists in ``doctor''. This bias gives rise to unfair job recommendations for the job seekers due to their registered gender information. 
%In another example, 
%Obviously, it is necessary to develop bias detection methods for word embeddings, in order to get rid of human biases in searches, recommendations or advertisements. 

%A word embedding represents each word as a d-dimensional word vector $\overrightarrow{w}\in\mathbb{R}^{d}$. In this case, the gender stereotype is disclosed by the word embeddings by: $\overrightarrow{\text{father}} - \overrightarrow{\text{mother}} \approx \overrightarrow{\text{doctor}} - \overrightarrow{\text{nurse}}$.

% Researchers have begun to look into the potential stereotypes in machine learning models
% Nevertheless, few research so far has been conducted to explore and solve the female/male gender stereotype problems in word embedding.


To alleviate gender stereotype in word embeddings, \newcite{kaiwei} propose a post-processing method that projects gender-neutral words to a subspace which is perpendicular to the gender dimension defined by a set of gender-definition words.\footnote{Gender-definition words are the words associated with gender by definition (e,g., mother, waitress); the remainder are gender-neutral words.} However, their approach has two limitations. First, the method is essentially a pipeline approach and requires the gender-neutral words to be identified by a classifier before employing the projection. If the classifier makes a mistake, the error will be propagated and affect the performance of the model. Second, 
% The word \textit{debias} was defined in that paper as the procedure of reducing the bias encoded in the learned representations of words.
%But this method fails to influence the word embedding training process. The post-processing fashion hurts the latent word similarities a lot. 
%However, all the gender neutral words are verified manually, which is time and cost consuming.
% Therefore, without the readjustment of embedding training model, 
%In addition, this method removes the latent gender information from the word embeddings, which is useful in several domains such as biomedical and social science~\cite{back2010gender,mcfadden1992study}.
%However, to get the gender-neutral word set, they need to train a classifier based on some seed words, which makes the neutralizing performance strongly depends on the quality of the classifier and the seed words.
%What is more,  
their method completely removes gender information from those words which are essential in some domains such as medicine and social science~\cite{back2010gender,mcfadden1992study}.

% at the cost of flexibility to use the word embeddings. 
% However, it will be better to keep it for further use while neutralizing the word embeddings.

To overcome these limitations, we propose a learning scheme, Gender-Neutral Global Vectors (GN-GloVe) for training word embedding models with protected attributes (e.g., gender) based on GloVe~\cite{glove}.\footnote{The code and data are released at \url{https://github.com/uclanlp/gn\_glove}} 
GN-GloVe represents protected attributes in certain dimensions while neutralizing the others during  training. As the information of the protected attribute is restricted in  certain dimensions, it can be removed from the embedding easily. By jointly identifying gender-neutral words while learning word vectors, 
GN-GloVe does not require a separate classifier to identify  gender-neutral words; therefore, the error propagation issue is eliminated. 
The proposed approach is generic and can be incorporated with other word embedding models and be
applied in reducing other societal stereotypes. 


%learns the embeddings with specific dimensions representing the desired attributes (e.g., gender). 
% Although certain attributes are presented separately, 
% Therefore, our approach still maintains the gender feature of each word. 
% In this work, we focus on the binary gender case, however the method can be generalized to reduce other stereotypes such as race or age bias.
% In this paper, we propose a new method to learn word representations invariant to gender stereotype, aiming at better applicability. 
%Comparing to the previous work, we do not need to train a separate classifier and can still keep the gender information which can be easily removed if needed.
%The generality of our approach enables word embedding
%models to handle different stereotypes in the learning phase, such as demographic bias and language stereotypes.
%Furthermore, representing specific features by predefined dimensions improves the interpretability of word embeddings.
%In this paper, we employ GloVe~\cite{glove} as the base embedding model.
%However, it is worth noting that the selection of the training model is not limited to GloVe.
%In other words, DeGloVe applies to other embedding methods.
% Inspired by the recent advancements of adversarial learning~\cite{GAN}, we formulate this representation learning problem as a simplified mini-max game among three players: a generator which captures global co-occurrence relations among words and then optimizes the embedding; a discriminator that identifies each word's gender type (female/male) from specific dimensions of its embedding vector; and an encoder which neutralizes the rest dimensions (debiased part) of the non-gender definition words which will be described in details in Section~\ref{sec:methodology}. 
% In another word, this generative adversarial method forces to transport the gender information from the whole vector to specific dimensions. In this mini-max game, we do not have the concerns of reaching an unstable Nash equilibrium as the original GAN does, because three players in our method will all converge in the final stage. 

Our contributions are summarized as follows: 1) To our best knowledge, GN-GloVe is the first method to learn word embeddings with protected attributes;
2) By capturing protected attributes in certain dimensions, our approach ameliorates the interpretability of word representations;
3) Qualitative and quantitative experiments demonstrate that GN-GloVe effectively isolates the protected attributes and preserves the word proximity. 

%Users can leverage different part of the word embeddings for various tasks according to the feature requirements. 

% models the stereotype in word embeddings as a generative adversarial problem. 

% The method can also worked as a new way to interpret word embeddings and based on which we can learn specific features representation.
% , if any related to the word.
% A more generic solution of interpreting word embeddings and learning feature representation is proposed. 
% Excessive experiments were conducted to evaluate the performance of the proposed model. In this work, we use GloVe~\cite{glove} as our fundamental  training method. Experimental results show that we can perverse the gender information using some specific dimensions while reducing the stereotype among other dimensions. And the performance on the word similarity and analogy tasks are almost the same comparing to the original model.
% Comparing to the hard debias method in~\cite{kaiwei}, we do not completely remove the gender information from the word embeddings. 


% The paper is organized as follows.
% Section~\ref{sec:related_work} gives a brief summary of different works working on the stereotype removing. 
% The major model framework is introduced in Section~\ref{sec:methodology} as well as the rationale behind it. 
% And section~\ref{sec:experiments} displays and explains the empirical results to support the proposed model.

